
Reinforcement learning from human feedback (RLHF) is a widely applied paradigm for aligning language model outputs with human preferences [Bai et al., 2022; Ouyang et al., 2022; Rafailov et al., 2023]. In this paradigm, human annotators compare pairs of model responses, and a reward model is trained to predict their preferences; a policy model is subsequently optimized against this reward. A known consequence of this process is the amplification of surface-level stylistic features that correlate with annotator approval, including structured enumeration, hedged assertions, and explicit reasoning chains [Stiennon et al., 2020; Rafailov et al., 2023].

A related question — whether prolonged human interaction with RLHF-optimized outputs in turn modifies human annotator preferences — has received comparatively little empirical attention. This constitutes one direction of bidirectional human-AI alignment, in which AI systems and human behavior mutually influence each other over time. Shen et al. [2024] synthesize over 400 papers across alignment, human-computer interaction, and cognitive science, finding that while AI-to-human effects such as persuasion and automation bias appear in multiple experimental studies, the reverse direction — systematic human behavioral adaptation toward AI-idiomatic norms — has not been empirically quantified in the alignment literature. Vishwarupe et al. [2026] audit 16 alignment benchmarks and find that none include mechanisms for verifying annotator behavioral change over time.

The present study was designed to test one specific instantiation of the human-to-AI adaptation hypothesis: whether annotators with prior exposure to deployed RLHF-trained systems show a measurably stronger preference for AI-idiomatic stylistic features relative to annotators without such exposure. To operationalize this hypothesis, the study introduces the AIFS construct — a composite count measure covering four stylistic feature classes — and applies a conditional logit regression with semantic cluster fixed effects to the Anthropic HH-RLHF dataset \citep{Baietal2022}, using the dataset's helpful-online versus helpful-base split as a proxy for annotator AI-exposure history.

The main results are as follows. First, AIFS features significantly predict annotator preference overall (β₁ = 0.0246, p < 0.001), establishing the construct's validity as a preference signal. Second, the interaction between AIFS features and annotator condition (online versus base) is not statistically significant (β₄ = −0.0016, OR = 0.9984, 95\% CI [0.9861, 1.0108], p = 0.796), and all four pre-registered gate thresholds are not met. Third, examination of the dataset documentation reveals that the split labels record which model generated the presented responses, not which annotators had prior AI exposure. The split is a data provenance label, not an annotator characteristic, and its use as an exposure proxy is not supported by the available metadata.

This paper makes the following contributions. It introduces and validates the AIFS construct as a preference predictor in a large-scale preference corpus. It documents and resolves a numerical instability issue — BFGS optimizer divergence on conditional logit problems with large numbers of small cluster fixed effects — through Newton-Raphson optimization. It provides an empirical demonstration that split-based adaptation proxies in the HH-RLHF dataset do not correspond to annotator exposure conditions, a distinction that has not previously been drawn with supporting evidence. It outlines the data infrastructure requirements that would be necessary for a valid empirical test of the bidirectional adaptation hypothesis.

The remainder of this paper is organized as follows. Section 2 reviews related work. Section 3 describes the method. Section 4 details the experimental setup. Section 5 presents results. Section 6 discusses interpretations and limitations. Section 7 concludes.

